% Complete Chapter 4 introduction, PDF 43 / printed 33.
% block: hilbert-complexes-heading page:43 type:chapter
\clearpage\setcounter{section}{3}
\section{{Hilbert complexes}}\label{chap:hilbert-complexes}

% block: hilbert-complexes-introduction page:43 type:ordinary-paragraph
{We have just seen that the differential operators entering the de Rham complex can be viewed as closed unbounded operators between $L^2$ spaces. In this way we obtain a structure which combines the homological algebraic features of a cochain complex, discussed in Chapter~\ref{chap:homological-algebra}, with the functional analytic features of closed unbounded operators on Hilbert space, discussed in Chapter~\ref{chap:unbounded-operators}. This structure is called a \emph{Hilbert complex} and was first promulgated by Brüning and Lesch \BookCite{24}. It is the structure underlying Hodge theory of Riemannian manifolds,}%
% block: hilbert-complex-history-footnote page:43 type:footnote
\footnote{{An indication of how natural is the notion of a Hilbert complex is that in the early drafts of \BookCite{13} my coauthors and I introduced this definition, before we discovered that it already occured in the literature \BookCite{24} with the very same name and definition.}}%
% block: hilbert-complexes-introduction-continuation page:43 type:ordinary-paragraph
{and, despite its simplicity, the main results of Hodge theory, for example, the identification of cohomology with harmonic forms and the Hodge decomposition, follow easily in this framework. Associated to each Hilbert complex is a graded map, or sequence of operators, called the abstract Hodge Laplacian of the complex. We explore several different but equivalent formulations of the Hodge Laplacian and demonstrate their well-posedness up to harmonic forms. The case of the de Rham complex on a three-dimensional domain is treated in detail, and the results of the chapter are summarized in Box~\ref{box:hilbert-complexes-summary} at the end of the chapter.}
